\documentclass[10pt,a4paper]{amsart}
\usepackage[margin=19mm]{geometry}
\usepackage{amsmath,amssymb,mathtools}
\usepackage[scr=rsfs,cal=euler]{mathalfa}
\usepackage{xfrac}
\usepackage[unicode,hidelinks,bookmarks=false]{hyperref}
\newcommand{\ie}{i.e.\ }
\newcommand{\cf}{cf.\ }
\newcommand{\resp}{respectively\ }
\newcommand{\Subsection}{\subsection}
\providecommand{\nobreakdash}{\nobreak\hbox{-}}
\setlength{\emergencystretch}{2em}
\multlinegap0pt
\usepackage{amssymb}
\usepackage[shortlabels]{enumitem}
\usepackage{mathtools}
\usepackage{bm}
\newtheorem{definition}{Definition}
\newtheorem{theorem}{Theorem}
\usepackage{cleveref}
\crefname{definition}{Definition}{Definitions}
\crefname{theorem}{Theorem}{Theorems}
\newcommand{\R}{\mathbb{R}}
\newcommand{\Z}{\mathbb Z}
\newcommand{\p}{\partial}
\title{Discrete Octonionic Dirac Operators on Bounded Lattices: Boundary Integral Theory and Scaling Limits}
\author{Guangbin Ren, Xin Zhao}
\date{Source: arXiv:2312.12777v2 (CC BY 4.0)}
\begin{document}
\maketitle

\noindent\emph{Source and licence.} Discrete Octonionic Dirac Operators on Bounded Lattices: Boundary Integral Theory and Scaling Limits. Version arXiv:2312.12777v2 (CC BY 4.0), distributed by arXiv under the Creative Commons Attribution 4.0 International licence, http://creativecommons.org/licenses/by/4.0/, as recorded in the arXiv licence metadata for this exact version and shown on the abstract page https://arxiv.org/abs/2312.12777v2. Full licence notice: \texttt{licenses/arxiv-2312.12777-CC-BY-4.0.txt}.\par\smallskip

\noindent\textit{Editorial scope.} Counted unit: the article's theorem theorem (source0.tex lines 1597--1606). The article's complete original proof of it is delivered with it (source0.tex lines 1608--1725, one \texttt{proof} environment of 118 lines). No mathematics is written, repaired or re-derived: the statement and the proof are the article's own lines, in the article's own notation, and no retained line is substituted or re-flowed.  Retained uncounted as the source-local context the proof needs: lines 1571--1591.  Nothing was omitted from any retained span, so the package carries no omission record.  No retained line contains a citation, so the wrapper carries no bibliography and no bibliography entry.  No retained line contains a figure environment, an image-inclusion command or drawing code, so no image is delivered and the package declares no asset dependency.  Editorial formatting only, in addition: an AMS \texttt{amsart} preamble that restates the article's own declarations and macros used by the retained lines, namely the environments \texttt{definition}, \texttt{theorem} on separate counters, together with the standard packages those lines need.  The retained environments print in this document as Definition 1, Theorem 1 in order of appearance, whereas the article prints the same items with the numbering of its own full document.  The complete native source of the article as distributed in the arXiv e-print for this exact version is bundled unchanged as \texttt{sources/151802/source0.tex}.
\begin{definition}\label{2023-09-11 10:00:54}
For a bounded domain $B$ within $\R^8$,	
	a collection of sets $B^h\subset \Z_h^8$ with $h>0$ is said to converge to $B$, expressed as
   \[
\lim_{h\to 0} B^h= B,
\] when the following conditions hold:
    \begin{equation}
        \begin{aligned}
            \lim_{h\to 0} \max_{\alpha\in \p B }\min_{\beta\in \p B^h} \Vert \alpha-\beta\Vert=0,\\
            \lim_{h\to  0} \max_{\alpha\in \p B^h }\min_{\beta\in \p B } \Vert \alpha-\beta\Vert=0,\\
            \lim_{h\to  0} \max_{\alpha\in \overline B}\min_{\beta\in B^h } \Vert \alpha-\beta\Vert=0,\\
            \lim_{h\to  0} \max_{\alpha\in B^h }\min_{\beta\in \overline B } \Vert \alpha-\beta\Vert=0.
        \end{aligned}
    \end{equation}

Furthermore, if $B^h$ converges to $B$ and
\[
V^h(B^h\setminus B) = O(h^2),
\]
then we say that the family $B^h$ has \emph{exterior excess of order $O(h^2)$}.
\end{definition}

\begin{theorem}
	For a bounded domain  $B$ in $\R^8$, the sets
  \[
B^h \equiv (B \cap \Z_h^8)^\circ
\] converge to $B$. Moreover,
  \[
V^h(B^h\setminus B)=0
\]
  for every $h>0$.
\end{theorem}

\begin{proof} Let us  confirm the four convergence identities from \Cref{2023-09-11 10:00:54}.

    \bigskip
    (i) First we verify the fourth identity. Since $B^h$ is contained in $B$, we have \begin{equation*}
        \min_{\beta\in \overline B } \Vert \alpha-\beta\Vert=0, \quad \forall \,\alpha\in B^h
    \end{equation*}
    so that  \[
\lim_{h\to 0} \max_{\alpha\in B^h }\min_{\beta\in \overline B } \Vert \alpha-\beta\Vert=0.
\]

    \bigskip

    (ii) To verify the third identity, we aim to show that given any $\varepsilon>0$, there exists a $\delta>0$ ensuring \[
\overline B \subset B^h + B(0, 2\varepsilon	)
\]   for any $h \leqslant \delta$.

    For any given 	$\varepsilon>0$ and $x \in \overline  B$ one can find $\delta > 0$ and $x^\prime \in  B$ such that \[
B(x^\prime,\delta ) \subset\subset B \cap B(x, \varepsilon	).
\]
    Consequently,
    \[
B(x^\prime,\delta) \cap \Z^8_h \subset  (B \cap \Z^8_h)^\circ  = B^h
\]
     for sufficiently small
  $h$, implying
     \[
B(x^\prime,\delta) \cap B^h = B(x^\prime,\delta) \cap \Z^8_h \neq \emptyset
\] for sufficiently small
     $h$.    From the fact \[
B(x,\varepsilon 	) \cap B^h \supset B(x^\prime,\delta) \cap B^h,
\] it follows that \begin{equation}
        \label{2023-09-11 16:15:35}
        B(x,\varepsilon 	) \cap B^h \neq \emptyset,\quad \forall h\ll 1.
    \end{equation}
    Due to the compactness of $B$, we can find a sequence $x_1,\cdots,x_n$ such that
    \[
\overline B\subset \bigcup\limits_{l=1}^n B(x_l,\varepsilon).
\] As we have proved for each $x_l$ there exists $\eta_l > 0$ such that \[
B(x_l,\varepsilon)\cap B^h  \neq \emptyset
\] for any $h<\eta_l.$ Hence for any $h<\min\limits_{1\leqslant l\leqslant n}\{\eta_l\},$ we have  \[
\overline B \subset B^h + B(0, 2\varepsilon	).
\]
    \bigskip

    (iii) Now we come to prove the second identity. It suffices to prove that \[
\min_{y\in \p B} \Vert x-y\Vert \leqslant 2h
\] for any  $x\in \p B^h$.

    {\bf Case 1:}  $x\in \p^+ B^h:= \p B^h\cap B^h.$

    \bigskip
    For any $y\in \p^+(B\cap \Z^8_h),$ according to the definition of $\p^+(B\cap \Z^8_h),$ we have $y \in B$ and there is a point $y^\prime \in \Z^8_h\setminus B$ such that $\Vert y^\prime-y \Vert = h.$ The fact that $y \in B$ and $y^\prime \notin  B$ leads to the conclusion that \[
\mathop{\min}\limits_{\beta \in \p B}\Vert y-\beta \Vert \leqslant h
\]
   for all  $y\in \p^+(B\cap \Z^8_h)$.






    For any $x\in \p^+ B^h,$ there exists $y\in N(x)$ such that $y\notin B^h$. Since $B^h$ is the discrete interior of $B\cap \Z^8_h,$ we have $y\in \p^+(B\cap \Z^8_h),$ which implies \[
\min_{\beta \in \p B}\Vert x-\beta \Vert \leqslant \Vert x-y\Vert +\min_{\beta\in \p B}\Vert y-\beta \Vert \leqslant 2h.
\]

    \bigskip
   {\bf  Case 2: } $x\in \p^-B^h:=\p B^h\setminus B^h.$

    \bigskip
    From the fact that $B^h=(B\cap \Z^8_h)^\circ,$ it follows that \[
x\in \p^+(B\cap \Z^8_h).
\] As shown above in case 1, we find \[
\min_{\beta\in \p B} \Vert x-\beta \Vert \leqslant h.
\]

    \bigskip

    (iv) At the final step, we come to verify the first identity \[
\lim_{h\to 0} \max_{\alpha\in \p B }\min_{\beta\in \p B^h} \Vert \alpha-\beta\Vert=0.
\] Since $\p B$ is compact, it is sufficient to prove that for any 	 $\varepsilon>0$  and $x \in\p B$, there exists $\delta>0$ such that \[
\min_{\beta \in \p B^h} \Vert x-\beta \Vert < \varepsilon
\] for any $h<\delta.$

    Indeed, for any given 	$\varepsilon>0$ and $x \in \p B$, as shown in \eqref{2023-09-11 16:15:35}, there exists $\eta>0 $ such that \[
B(x,\varepsilon)\cap B^h\neq \emptyset
\] for any $h<\eta.$

    On the other hand, because $B(x,\varepsilon 	)\setminus B$ is open, there exists $\eta^\prime > 0$ such that \[
B(x,\varepsilon)\cap \Z^8_h\setminus \overline B \neq \emptyset
\]
   for all  $h<\eta^\prime$.

    When $h<\min\{\eta,\eta^\prime\}$, choose
    \[
        \alpha \in B^h\cap B(x,\varepsilon), \qquad \beta \in (\Z_h^8\setminus B^h)\cap B(x,\varepsilon).
    \]
    Because the lattice set \(B(x,\varepsilon)\cap \Z_h^8\) is connected by nearest-neighbor steps, there exists a sequence
    \[
        \{x_l\}_{l=1}^n \subset \Z_h^8 \cap B(x,\varepsilon)
    \]
    such that
    \[
        x_1=\alpha,\qquad x_n=\beta,\qquad \Vert x_l-x_{l+1}\Vert=h
    \]
    for all \(1\le l\le n-1\). Since \(x_1\in B^h\) and \(x_n\notin B^h\), there exists an index \(l_0\) with
    \[
        x_{l_0}\in B^h,\qquad x_{l_0+1}\notin B^h.
    \]
    By the definition of the discrete boundary, both \(x_{l_0}\) and \(x_{l_0+1}\) belong to \(\p B^h\). In particular,
    \[
        \min_{\gamma\in \p B^h}\Vert x-\gamma\Vert<\varepsilon.
    \]
    This proves the first convergence identity. Since \(B^h\subset B\), we also have \(V^h(B^h\setminus B)=0\), and the proof is complete.




\end{proof}

\end{document}
